\scanpage{33}
\mspage{13}

$W(k)$ des vecteurs de Witt sur $k$, anneau de val. discrète dont le corps des fractions $W(k)[\frac1p]\supset\mathbb Q_p$ est un corps de car. \uline{nulle}.

Ceci devrait être le cas encore pour un corps de base \inserted{parfait} quelconque, non pas néc. alg. clos \ldots

\S 3 \uline{Revue des essais de définition d'une théorie \inserted{cohomologique} $p$-adique en car. $p>0$.}

\uline{Retour sur De Rham via Washnitzer Monsky.}

\inserted{\uline{Lien avec les stratifications de Gauss-Manin.}}

1) Utilisation des $p$-groupes éventuellement radiciels, \deleted{pour} comme coefficients pour la cohomologie (\unsure{pour|sur} la topologie plate).

Heuristiquement, cela marche magnifiquement en dim. 1, et \deleted{\illegible{3} \unsure{pour}} \inserted{on entrevoit une variante considérablement améliorée, de nature géométrique,} \deleted{\unsure{de la th}} en termes \unsure{du théorème|des théorèmes} de dualité en cohomologie plate, du corps de classes global géométrique (\unsure{constr.} sur un

\margin{32}
